\chapter{What Zeno Does Not Have Yet}\label{app-limits}

Before version 0.1, every decision in \verb|docs/decisions.md| was read
again, and thirty were answered "a limit of 0.1": the thing is left out
of this version, and a program that asks for it is told so. Later
versions have taken some of them off the list, and what follows is
what is still left out. Each can be added later without changing what any program written now means.
The numbers are the decisions, for anyone who wants the argument.

\section{On a board}

\begin{description}
\item[Inputs.] A pin can be read (\verb|input_pin|, 45), a serial
  port used (\verb|serial|, 46), a voltage read
  (\verb|analog_pin|, 47) and a PWM signal made (\verb|pwm_pin|, 48),
  but nothing presses a button or feeds a
  voltage under \verb|zeno run| (45.5), and none of it has been tried
  on a real board. A PWM pin cannot be asked its duty or given a new
  frequency (48.6). A board keeps sixteen tickers and none is given
  back (\verb|ticker|, 49.5). An interrupt counts edges
  (\verb|interrupt|, 50), and no function of the program runs inside
  one (50.3); nor can a program wait for the first of several.
\end{description}

\section{Numbers}

\begin{description}
\item[A \texttt{u64} in a few places.] A \verb|u64| is not yet a key
  of a dict, an index, a count for \verb|range|, or something
  \verb|json.write| can write; \verb|int( )| makes one an \verb|int|
  (42.10).
\item[A list of a fixed length takes a block.] A \texttt{list[u8, 64]}
  in a build with no collector is a block from the pool, as any list is,
  though its size is known (36.5).
\item[Units of a program's own.] Units make units only from a closed
  list (\verb|V * A| is \verb|W|, \verb|1 / 2 ms| is \verb|Hz|), and a
  program cannot name a unit of its own, or use charge (\verb|mAh|) or
  speed yet (43.6).
\end{description}

\section{Types}

\begin{description}
\item[Unions.] Only \verb|or none| can follow a type: \verb|int or str|
  is refused (1.13). An enum that carries data does the same work.
\item[One bound for a type.] A type parameter asks for at most one
  interface, such as \verb|[T is Ordered]|, not two at once, and a
  class cannot be \verb|Ordered| until it can say what \verb|<| means
  for its objects (58.3).
\item[A generic as a value.] A generic function, or a case of a
  generic enum, cannot be kept as a value, a method cannot have type
  parameters of its own, and the types cannot be written at a call
  (3.16e, 58.1).
\item[Several values in a \texttt{for}] (\verb|for q, r in pairs()|),
  and a default for one of several results (57.6).
\item[A resource lives only in a \texttt{let} or a parameter:] not in a
  field, a list, a dict, \verb|or none|, an enum or a type parameter,
  and a \verb|release| cannot raise (59.2).
\end{description}

\section{The library}

\begin{description}
\item[A file's name on Windows.] In a program built through C or LLVM
  on Windows, a file's name goes through the system's code page, as
  \verb|fopen| takes it, so a name with letters that page lacks cannot
  be opened there; \verb|zeno run| takes any name (59.20).
\item[Files in pieces.] A reader cannot go back or say where in the
  file it is, does not wait for a file to grow, and reads only UTF-8;
  each \verb|write| goes to the system at once (59.6).
\item[The methods of text and lists written in Zeno.] A part of the
  library is written in Zeno itself: \verb|math.hypot|, the hyperbolic
  functions, \verb|math.gcd|, \verb|math.clamp| and the module
  \verb|random| (33.1). The methods of text and lists are still built
  into the compiler (6.10). And \verb|random| has no seed from the
  machine: a program gets the same numbers on every run until it
  calls \verb|random.seed| (33.8).
\end{description}

\section{Libraries of other languages}

\begin{description}
\item[A function of Zeno cannot be given to another language]: not as
  a callback of C, a block of Ruby, a lambda of LINQ or a sort key of
  Python (26.12, 44.9).
\item[The seven that run beside Zeno] (\verb|import python| and its
  like) are for \verb|zeno run| and \verb|zeno test|. Every build
  refuses them (26.12).
\item[From C:] a struct given whole, a function with \verb|...|, the
  address of a variable, and declarations read from a header by
  themselves are not there (44.9). \verb|--llvm| and \verb|--target|
  refuse \verb|import c|, and it has not been run with a
  cross-compiler's own libraries or on a board.
\item[An \texttt{inline} block] is of C, Python or JavaScript, and
  not of Ruby, PHP, Lua, C\#, C++ or assembly by itself. A block of C
  cannot name a library it needs, and Zeno sees its functions and not
  its constants. \verb|--llvm| and \verb|--target| refuse a block,
  and every build refuses one of Python or JavaScript (61.11, 61.19).
\item[A call the C compiler can work out.] A built program's
  \verb|m.cbrt(27.0)| is worked out by GCC before the program runs,
  rounded exactly, where \verb|zeno run| asks the library: the two can
  differ in the last digit.
\end{description}

\section{Builds}

\begin{description}
\item[One file at a time.] A build compiles the whole program at once
  (7.11).
\item[A board by its name.] \verb|zeno build --board| takes the
  folder of a board, which today is in a copy of Zeno's repository or
  written by you; it does not load the program onto the board, and a
  build for a board has no collector (51.9). The two boards that come
  with Zeno are one emulated chip.
\item[\texttt{--memory} on a board that prints by itself.] The
  report is written as the C library ends a program, and such a board
  has no C library under the program (52.8).
\item[A NaN with decimals in a minimal build for a computer.] There
  \verb|{x:.2f}| is \verb|printf|'s own, so a float that is not a
  number is shown as the C library writes it, \verb|nan|. Every other
  build prints \verb|NaN|, as \verb|zeno run| does (53.8).
\item[A compiler that is neither GCC nor Clang.] It gets the checks
  in plain C (chapter 55), and tcc builds every example so. No program
  has been built by IAR's, Keil's or a chip maker's compiler. A
  firmware build with a collector cannot be built by such a compiler,
  and its C says so (55.5).
\item[\texttt{setjmp} in an ordinary build.] An ordinary build uses it
  for \verb|try|, and a firmware build does not use it at all (5.11l,
  12.8).
\end{description}

\section{Firmware and no collector}

\begin{description}
\item[What a build with no collector refuses by design.] Replacing
  what a field holds through a parameter that only looks, taking an
  item out of a list a \verb|for| walks (5.15e), keeping an item that
  \verb|[ ]| only looks at, and a kept function that uses a \verb|var|
  from outside it. Each is refused with the reason, and the table of
  \verb|docs/no-collector.md| lists them all (56.17).
\item[Swapping two items a list owns,] \verb|xs[i], xs[j] = xs[j], xs[i]|
  on a list of lists or objects with no collector: each item is only
  looked at there (57.6).
\end{description}

\section{Interfaces}

\begin{description}
\item[Functions of an interface] have no type parameters and cannot
  be private (10.5), and a default for one of their parameters is a
  plain value (58.20).
\item[An interface that is another interface] (10.10, 58.23).
\end{description}

\section{The web}

\begin{description}
\item[Only what the module has.] No images drawn on a canvas, no
  curves, no touch with more than one finger, and no address bar of
  its own (16.13).
\item[An interpreter in a page.] Every page carries \verb|zeno.wasm|
  and runs at the speed of \verb|zeno run|. A page compiled through C
  or LLVM, small and fast, is left for later (14.1).
\item[A simulated page reads plain selectors only,] and knows no CSS
  (16.9).
\item[The playground] runs one file, and has no colours in its box and
  no Stop button (16.12).
\item[JSON] is read into dicts and lists, not straight into a class,
  and is not in compiled programs (18.5).
\end{description}

\section{Tools}

\begin{description}
\item[Going to a field.] The language server shows the type of a
  field, but cannot yet jump to where it was declared (17.4).
\item[\texttt{assert} is only for tests] (10.21).
\end{description}

\section{Not limits}

Three things are not on this list, because they are decided: Zeno has
no inheritance, no second way to write a function, and no second name
for a module. The reasons are in \hyperref[ch-design]{Why Zeno Is as It
Is}.
